\title{The Era of 1-bit LLMs: All Large Language Models are in 1.58 Bits}
\begin{equation}
    \mathrm{RoundClip}(x, a, b) = \max(a, \min(b, \mathrm{round}(x))),
\end{equation}

We extended our analysis to larger model scales, specifically 7B, 13B, and 70B parameters, and assessed the associated computational costs. As presented in Figure~\ref{fig:latency-memory-energy}, both latency and memory requirements exhibit clear scaling patterns; notably, the performance improvement grows with model size. For instance, \text{BitNet b1.58} 70B demonstrates a 4.1-fold speed-up in comparison to the \text{LLaMA LLM} reference model. This acceleration is primarily attributed to the increased computational burden of the \emph{nn.Linear} operation in larger models. Similarly, memory utilization follows an analogous scaling behavior, since the embeddings retain full-precision representation and their share of total memory usage becomes less significant as models increase in size. Both performance metrics—latency and memory—were measured using a 2-bit kernel, suggesting further optimizations remain possible to lower computational cost.

\paragraph{Throughput}

For activation quantization, we use the same technique as BitNet; however, unlike the original approach, we do not rescale activations into $[0, Q_b]$ prior to applying non-linear transformations. Instead, activations are rescaled for each token to $[-Q_b, Q_b]$, thereby eliminating the necessity for zero-point quantization. This alternative is simpler and more straightforward for practical implementation and system-level tuning, with only minimal impact on model accuracy observed in our evaluations.

\begin{itemize}
    \item BitNet b1.58 with 13B parameters demonstrates superior latency, memory consumption, and energy efficiency relative to a 3B FP16 LLM.
    \item BitNet b1.58 at 30B parameters surpasses a 7B FP16 LLM across latency, memory, and energy metrics.
    \item BitNet b1.58 70B outperforms a 13B FP16 LLM with respect to latency, memory footprint, and energy use.
\end{itemize}

\textbf{New Hardware for 1-bit LLMs}

\section{The Era of 1-bit LLMs}

\paragraph{Training with 2T Tokens}

\textbf{LLMs on Edge and Mobile}

The rapid expansion of Large Language Models (LLMs) in recent years has greatly advanced the capabilities and versatility of artificial intelligence in natural language processing. Despite these achievements, the trend toward larger models introduces deployment obstacles and amplifies concerns about financial and environmental sustainability due to escalating energy requirements.

\textbf{\text{BitNet b1.58} enables a novel scaling law that connects model performance with inference efficiency.} Drawing on the data presented in Figures~\ref{fig:latency-memory-energy} and~\ref{fig:energy}, we can establish approximate equivalencies between models of different sizes when utilizing 1.58-bit and 16-bit precision:

\begin{equation}
    \widetilde{W} = \mathrm{RoundClip}(\frac{W}{\gamma + \epsilon}, -1, 1),
\end{equation}


\begin{abstract}
Emerging work, notably BitNet~\cite{bitnet}, is ushering in a paradigm shift toward 1-bit Large Language Models (LLMs). This paper presents a 1-bit LLM variant named \textbf{\text{BitNet b1.58}}, in which each individual model parameter takes on a ternary value from the set \{-1, 0, 1\}. This architecture achieves performance parity—both in terms of perplexity and downstream task effectiveness—with full-precision Transformer LLMs (utilizing FP16 or BF16), given the same parameter count and training data, while greatly reducing resource costs in latency, memory usage, throughput, and energy expenditure. More fundamentally, the 1.58-bit LLM establishes a novel scaling principle and methodology for training advanced LLMs that unite high efficiency with superior performance. Additionally, it paves the way for a new computational model and motivates the development of bespoke hardware architectures tuned for 1-bit LLMs.
\end{abstract}

The perplexity and computational cost for both \text{BitNet b1.58} and \text{LLaMA LLM} are presented in Table~\ref{tab:ppl}. The data indicate that at a 3B parameter count, \text{BitNet b1.58} achieves perplexity comparable to that of its full-precision \text{LLaMA LLM} counterpart, all while being 2.71$\times$ faster and requiring only 28\% ($\approx$1/3.55) of the GPU memory. Notably, with a model size of 3.9B parameters, \text{BitNet b1.58} outpaces the 3B \text{LLaMA LLM} by a factor of 2.4 in speed, has 3.32$\times$ lower memory consumption, and delivers significantly enhanced performance.

Token count plays a significant role in the effectiveness of LLMs. To evaluate how well \text{BitNet b1.58} scales with token volume, we trained it on 2T tokens, adhering to the data curation methodology of StableLM-3B~\cite{StableLM-3B-4E1T}, which currently represents the leading open-source 3B model. We compared both models using a benchmark set comprising Winogrande~\cite{winoGrande}, PIQA~\cite{piqa}, SciQ~\cite{sciq}, LAMBADA~\cite{lambada}, and ARC-easy~\cite{arc}. Zero-shot performance was calculated and summarized in Table~\ref{tab:2t}. For metrics reported as both accuracy and normalized accuracy, their mean was used. StableLM 3B’s 2T token results are sourced from its technical documentation. The findings demonstrate that \text{BitNet b1.58} consistently achieves the highest performance across all tasks, providing evidence that 1.58-bit LLMs possess strong generalization ability.

Detailed zero-shot evaluation results are included in Table~\ref{tab:zero-shot}. The experimental setup adhered to the procedure established by \emph{lm-evaluation-harness}\footnote{\url{https://github.com/EleutherAI/lm-evaluation-harness}}. Analysis of the findings reveals a shrinking disparity in performance between \text{BitNet b1.58} and \text{LLaMA LLM} as model scale increases. Critically, from the 3B parameter configuration onwards, \text{BitNet b1.58} achieves parity with the full-precision baseline. Echoing the perplexity trends, \text{BitNet b1.58} at 3.9B not only surpasses \text{LLaMA LLM} 3B in end-task accuracy, but does so with reduced latency and memory requirements, establishing \text{BitNet b1.58} as a Pareto superior alternative to leading LLM architectures.

\section{Results}

\section{BitNet b1.58}

\textbf{Native Support of Long Sequence in LLMs}

\paragraph{Quantization Function.} To limit weight values to -1, 0, or +1, we employ an \emph{absmean} quantization approach. This method involves normalizing the weight matrix using the mean absolute value of its elements and then rounding each element to its closest value among \{-1, 0, +1\}:

As LLMs evolve, supporting longer input sequences has become increasingly important. A key bottleneck is the substantial memory required to store KV caches during inference. \text{BitNet b1.58} offers promising progress by halving activation storage from 16 bits to 8 bits, thereby permitting twice the context length within the same hardware budget. Prospective advancements may allow further, lossless reductions of activation sizes to 4 bits or even less, specifically for 1.58-bit LLMs—an avenue reserved for future exploration.

We additionally approximated the energy consumption arising from arithmetic operations for both \text{BitNet b1.58} and \text{LLaMA LLM}. The evaluation centers on matrix multiplication, as it constitutes the primary contributor to total computation cost in LLMs. Figure~\ref{fig:energy} breaks down the energy composition, revealing that the majority of computations in \text{BitNet b1.58} employ INT8 addition, while \text{LLaMA LLM} largely relies on both FP16 addition and multiplication. Drawing on the energy consumption model provided in~\cite{energycost, pokebnn}, it was observed that \text{BitNet b1.58} reduces arithmetic energy usage for matrix multiplication by a factor of 71.4 on 7nm process technology. We also report holistic energy usage for end-to-end inference with 512 token inputs. The findings indicate that the energy efficiency advantage of \text{BitNet b1.58} increases with model scale relative to the FP16 \text{LLaMA LLM} baseline. This trend can be ascribed to the rising proportion of computation attributable to \emph{nn.Linear} layers in larger models, while the influence of other components diminishes.

\section{Discussion and Future Work}

Mixture-of-Experts (MoE) approaches have demonstrated their efficiency as a scalable solution for deploying LLMs by reducing computational FLOPs. However, their large memory requirements and the associated inter-chip communication overhead still pose significant obstacles to widespread use. Employing 1.58-bit LLMs can mitigate these issues: with a smaller memory demand, fewer devices are needed to host MoE architectures. Additionally, the bandwidth required for inter-device activation transfer is substantially lowered. If it becomes feasible to accommodate the whole MoE model on a single chip, this communication cost could be eliminated entirely.

Here, we present a notable 1-bit LLM variant named \textbf{\text{BitNet b1.58}}, in which each parameter takes on one of three values from the set \{-1, 0, 1\}. The inclusion of the zero value augments the original BitNet design, raising its precision to 1.58 bits in the binary system. \text{BitNet b1.58} preserves all the key advantages of its 1-bit predecessor: it employs a computation strategy that eliminates nearly all multiplication operations during matrix multiplication and can be extensively optimized. Energy consumption remains comparable to the original BitNet, while it also outperforms FP16 LLM baselines in memory efficiency, throughput, and latency. Moreover, \text{BitNet b1.58} introduces two further benefits. First, its expressiveness is enhanced by explicit feature filtering, enabled through the allowance of zero-valued parameters—this boosts the representational power of 1-bit LLMs. Second, empirical results demonstrate that \text{BitNet b1.58} achieves parity with FP16 models in both perplexity and end-task results, beginning at the 3B parameter scale, when matched for architectural and training configurations (including model size and the amount of training data).

1.58-bit LLMs promise substantial benefits for deploying language models on edge and mobile platforms, which often face stringent constraints in memory and computation. The efficiency gains from the reduced bitwidth enable broader deployment scenarios previously unattainable for LLMs, unlocking new functionalities and use cases for such devices. Furthermore, as these models are inherently more efficient in terms of both memory and power, they align well with CPU-dominant edge and mobile systems. Thus, \text{BitNet b1.58} can be executed efficiently on resource-constrained hardware, further extending the reach of LLM technology.

\textbf{1-bit Mixture-of-Experts (MoE) LLMs}

LLaMA~\cite{llama, llama2} has emerged as a standard architectural reference for open-source large language models. To align with the broader open-source ecosystem, \text{BitNet b1.58} incorporates LLaMA-style design choices, including RMSNorm~\cite{rmsnorm}, SwiGLU~\cite{swiglu}, rotary embedding~\cite{rope}, and the removal of all bias terms. Through this design, \text{BitNet b1.58} is readily compatible with commonly used open-source frameworks (such as Huggingface, vLLM~\cite{vllm}, and llama.cpp\footnote{\url{https://github.com/ggerganov/llama.cpp}}) with minimal integration work.

A viable solution has been the adoption of post-training quantization, which allows LLMs to operate with fewer bits during inference~\cite{smoothquant, gptq, quip, quip_sharp}. By lowering the bitwidth of model weights and activations, this approach leads to substantial savings in both memory and computation. While quantization strategies have evolved from 16 bits to as little as 4 bits~\cite{gptq, awq}, post-training quantization remains fundamentally limited, even if it has seen widespread deployment in industrial LLM systems.

\paragraph{Energy}

Recent advances in 1-bit model architectures, exemplified by BitNet~\cite{bitnet}, highlight a compelling path toward reducing the computational demands of LLMs without sacrificing efficacy. Standard LLMs employ 16-bit floating-point representations (FP16 or BF16), with matrix multiplication constituting the primary computational bottleneck. This results in significant costs due to intensive floating-point addition and multiplication operations. In contrast, BitNet's matrix multiplications rely exclusively on integer addition, yielding dramatic reductions in energy consumption for LLM workloads. Since many modern hardware platforms are constrained by power limitations, these energy efficiencies often correspond to improvements in processing speed.

\paragraph{LLaMA-alike Components.}

\paragraph{Memory and Latency}

\text{BitNet b1.58} is derived from the BitNet architecture, a Transformer variant in which traditional \emph{nn.Linear} layers are substituted with \emph{BitLinear}. This model is trained from the beginning using 1.58-bit precision for weights and 8-bit quantization for activations. In contrast to the initial BitNet design, it implements several updates, which are outlined below.

Our study involved benchmarking \text{BitNet b1.58} against our own FP16 implementation of the \text{LLaMA LLM} at multiple model scales. For a consistent evaluation, both models underwent pre-training on the RedPajama dataset~\cite{redpajama}, with exposure to 100 billion tokens. We assessed zero-shot generalization on an assortment of language understanding tasks: ARC-Easy~\cite{arc}, ARC-Challenge~\cite{arc}, Hellaswag~\cite{hellaswag}, Winogrande~\cite{winoGrande}, PIQA~\cite{piqa}, OpenbookQA~\cite{openbookqa}, and BoolQ~\cite{boolq}. Validation perplexities were calculated on both WikiText2~\cite{wikitext2} and C4~\cite{c4} datasets.

\begin{equation}
    \gamma = \frac{1}{nm}\sum_{ij} |W_{ij}|.
\end{equation}

Recent innovations exemplified by Groq\footnote{\url{https://groq.com/}} highlight the opportunities in developing purpose-built hardware, such as LPUs, tailored to the specific needs of LLMs. Looking ahead, there is significant potential—and an open call—to engineer new computing architectures and systems designed explicitly for the requirements of 1-bit LLMs, leveraging the computation model introduced with BitNet~\cite{bitnet}.

Beyond computation, the movement of model parameters from DRAM to on-chip accelerator memory (such as SRAM) is another costly aspect of inference. While increasing SRAM size has been explored to enhance throughput, this solution is often prohibitively expensive relative to DRAM. In contrast to full-precision models, 1-bit LLMs present a much smaller memory footprint with respect to both storage capacity and bandwidth requirements. This lower demand results in shorter and less resource-intensive DRAM-to-SRAM transfers, which accelerates and streamlines inference tasks.

We assessed the throughput of \text{BitNet b1.58} and \text{LLaMA LLM} with 70B parameters, running both models on a pair of 80GB A100 GPUs. To accommodate the size of \text{LLaMA LLM} 70B, we employed pipeline parallelism~\cite{gpipe}. For both models, batch size was scaled up incrementally until the maximum GPU memory was reached, with a fixed sequence length of 512. According to Table~\ref{tab:throughput}, \text{BitNet b1.58} 70B is able to handle a batch size up to 11 times larger than that of \text{LLaMA LLM}, yielding an 8.9-fold increase in throughput.

In addition, we systematically compared the GPU memory usage and inference latency between the \text{LLaMA LLM} and \text{BitNet b1.58}. Measurements were performed using the FasterTransformer\footnote{\url{https://github.com/NVIDIA/FasterTransformer}} framework, which is designed to optimize large language model inference on GPU hardware. For \text{BitNet b1.58}, we also incorporated the 2-bit kernel provided by Ladder~\cite{ladder}. Our primary metric for inference efficiency was the time required to generate each output token.